% === START SEGMENT 3 ===

%% ============================================================
%% SECTION V: EXPERIMENTAL SETUP
%% ============================================================
\section{Experimental Setup}
\label{sec:experimental_setup}

\subsection{Datasets}
We evaluate RobustIDPS on five widely adopted intrusion detection benchmarks:

\begin{enumerate}
    \item \textbf{CICIDS2017}~\cite{Sharafaldin2018}: 2,830,743 flows with benign traffic and 14 attack types generated over five days using CICFlowMeter.
    \item \textbf{NSL-KDD}~\cite{Tavallaee2009}: A refined version of KDD'99 containing 125,973 training and 22,544 test records spanning four attack categories.
    \item \textbf{CIC-IoT-2023}~\cite{Neto2023}: A recent IoT-specific dataset with 33 attack types across heterogeneous device profiles.
    \item \textbf{UNSW-NB15}~\cite{Moustafa2015}: 2,540,044 records encompassing nine attack categories with contemporary attack signatures.
    \item \textbf{CSE-CIC-IDS2018}~\cite{CSEIDS2018}: An extended multi-day capture with seven attack scenarios and updated traffic profiles.
\end{enumerate}

Across all datasets, we extract 83 flow-level features via CICFlowMeter and unify the label space into 34 attack classes plus benign, enabling cross-dataset evaluation.

\subsection{Evaluation Metrics}
We report Accuracy, macro and weighted F1-score, Precision, Recall, False Positive Rate~(FPR), and AUC-ROC. For continual learning we additionally measure Backward Transfer~(BWT) and Forward Transfer~(FWT)~\cite{Lopez-Paz2017}. Adversarial robustness is quantified via robust accuracy under attack.

\subsection{Implementation Details}
The system is implemented in PyTorch~1.13 with a FastAPI back-end and React front-end, deployed at \texttt{robustidps.ai}. All models are trained using the Adam optimizer~\cite{Kingma2015} with learning rate $10^{-4}$, batch size 256, and 50~epochs with early stopping (patience~=~5). Experiments run on four NVIDIA~A100 GPUs.

\subsection{Baselines}
We compare against: Random Forest~(RF), XGBoost~\cite{Chen2016}, standalone CNN, standalone LSTM, and a standard MLP ensemble without continual learning or adversarial training.

%% ============================================================
%% SECTION VI: RESULTS AND ANALYSIS
%% ============================================================
\section{Results and Analysis}
\label{sec:results}

\subsection{Detection Performance}

Table~\ref{tab:detection} summarizes classification performance across all four primary benchmarks.

\begin{table}[!t]
\centering
\caption{Detection Performance Across Benchmark Datasets}
\label{tab:detection}
\renewcommand{\arraystretch}{1.15}
\begin{tabular}{@{}lcccc@{}}
\toprule
\multirow{2}{*}{\textbf{Model}} & \textbf{CICIDS2017} & \textbf{NSL-KDD} & \textbf{CIC-IoT} & \textbf{UNSW-NB15} \\
 & Acc / F1 & Acc / F1 & Acc / F1 & Acc / F1 \\
\midrule
Random Forest   & 97.1 / 95.3 & 95.8 / 94.1 & 94.6 / 93.2 & 93.4 / 91.8 \\
XGBoost         & 97.5 / 95.9 & 96.2 / 94.8 & 95.1 / 93.8 & 94.0 / 92.5 \\
CNN             & 97.8 / 96.4 & 96.0 / 94.5 & 95.5 / 94.3 & 94.3 / 93.0 \\
LSTM            & 97.6 / 96.1 & 96.4 / 95.0 & 95.3 / 94.0 & 94.1 / 92.7 \\
MLP Ensemble    & 98.3 / 97.2 & 97.0 / 95.8 & 96.2 / 95.1 & 95.1 / 93.9 \\
\midrule
\textbf{RobustIDPS} & \textbf{99.2} / \textbf{98.7} & \textbf{98.1} / \textbf{97.4} & \textbf{97.8} / \textbf{97.1} & \textbf{96.9} / \textbf{96.2} \\
\bottomrule
\end{tabular}
\end{table}

RobustIDPS consistently outperforms all baselines, achieving 99.2\% accuracy and 98.7\% macro-F1 on CICIDS2017. The gains are most pronounced on the challenging CIC-IoT-2023 and UNSW-NB15 datasets, where the heterogeneous ensemble and adversarial training provide complementary benefits. The FPR remains below 0.3\% on all benchmarks, which is critical for operational deployment where false alarms incur significant analyst workload.

\subsection{Continual Learning}

We partition the unified 34-class label space into five sequential tasks $\mathcal{T}_1, \ldots, \mathcal{T}_5$ and measure how well the model retains prior knowledge while adapting to new attack types.

\begin{table}[!t]
\centering
\caption{Continual Learning Performance Across Sequential Tasks}
\label{tab:continual}
\renewcommand{\arraystretch}{1.15}
\begin{tabular}{@{}lccccc@{}}
\toprule
\textbf{Metric} & $\mathcal{T}_1$ & $\mathcal{T}_2$ & $\mathcal{T}_3$ & $\mathcal{T}_4$ & $\mathcal{T}_5$ \\
\midrule
Acc after task (\%) & 99.1 & 98.6 & 98.2 & 97.9 & 97.5 \\
BWT (\%)            & --   & $-$0.3 & $-$0.5 & $-$0.7 & $-$0.9 \\
FWT (\%)            & --   & +1.2 & +1.0 & +0.8 & +0.7 \\
\midrule
\multicolumn{6}{@{}l}{\textit{Without EWC (fine-tuning only):}} \\
Acc after task (\%) & 99.1 & 96.2 & 93.5 & 90.1 & 86.8 \\
BWT (\%)            & --   & $-$2.8 & $-$4.1 & $-$5.6 & $-$7.3 \\
\bottomrule
\end{tabular}
\end{table}

% Placeholder for forgetting curves figure
\begin{figure}[!t]
\centering
\fbox{\parbox{0.9\columnwidth}{\centering\vspace{2em}%
\textbf{[Figure Placeholder]}\\Forgetting curves: accuracy on $\mathcal{T}_1$ as subsequent tasks are learned, comparing EWC+replay vs.\ fine-tuning.\vspace{2em}}}
\caption{Accuracy on initial task $\mathcal{T}_1$ as new tasks are sequentially introduced.}
\label{fig:forgetting}
\end{figure}

The EWC-augmented replay mechanism limits backward transfer degradation to $-$0.9\% after all five tasks, compared to $-$7.3\% for naive fine-tuning (Table~\ref{tab:continual}). Forward transfer remains positive throughout, indicating that learned representations generalize to novel attack families.

\subsection{RL Response Agent}

Table~\ref{tab:rl_response} evaluates the Constrained PPO response agent on action selection accuracy, mean response latency, and safety constraint violations.

\begin{table}[!t]
\centering
\caption{RL Response Agent Performance}
\label{tab:rl_response}
\renewcommand{\arraystretch}{1.15}
\begin{tabular}{@{}lc@{}}
\toprule
\textbf{Metric} & \textbf{Value} \\
\midrule
Action selection accuracy (\%)  & 96.8 \\
Mean response time (ms)         & 12.4 \\
Safety constraint violations (\%) & 0.07 \\
Cumulative reward (converged)   & 847.3 \\
Episodes to convergence         & 1,240 \\
\bottomrule
\end{tabular}
\end{table}

% Placeholder for reward convergence figure
\begin{figure}[!t]
\centering
\fbox{\parbox{0.9\columnwidth}{\centering\vspace{2em}%
\textbf{[Figure Placeholder]}\\Reward convergence of Constrained PPO over training episodes.\vspace{2em}}}
\caption{Cumulative reward convergence of the RL response agent during training.}
\label{fig:reward_convergence}
\end{figure}

The agent achieves 96.8\% action selection accuracy with a response latency of 12.4\,ms, well within the real-time requirements of operational IDS deployments. Critically, safety constraint violations remain at 0.07\%, satisfying the $<$0.1\% threshold enforced by the CPO-inspired cost constraint.

\subsection{Adversarial Robustness}

Table~\ref{tab:adversarial} reports clean accuracy, accuracy under attack, and robust accuracy after adversarial training (AT) for each of the six attack methods.

\begin{table}[!t]
\centering
\caption{Adversarial Robustness Under Six Attack Methods}
\label{tab:adversarial}
\renewcommand{\arraystretch}{1.15}
\begin{tabular}{@{}llccc@{}}
\toprule
\textbf{Attack} & $\boldsymbol{\varepsilon / \sigma}$ & \textbf{Clean} & \textbf{Attacked} & \textbf{Robust (AT)} \\
\midrule
FGSM~\cite{Goodfellow2015}   & $\varepsilon{=}0.05$ & 99.2 & 84.3 & 96.8 \\
PGD-20~\cite{Madry2018}      & $\varepsilon{=}0.03$ & 99.2 & 78.6 & 95.4 \\
C\&W~\cite{Carlini2017}      & $c{=}10$             & 99.2 & 81.2 & 96.1 \\
DeepFool~\cite{Moosavi2016}  & --                   & 99.2 & 82.7 & 96.5 \\
AutoAttack~\cite{Croce2020}  & $\varepsilon{=}0.03$ & 99.2 & 76.1 & 94.8 \\
GAN-based~\cite{Lin2022}     & $\sigma{=}0.1$       & 99.2 & 79.4 & 95.2 \\
\bottomrule
\end{tabular}
\end{table}

Without adversarial training, even moderate perturbation budgets degrade accuracy by 15--23 percentage points. Our multi-attack adversarial training regime recovers the majority of this loss, achieving robust accuracies of 94.8--96.8\% across all six threat models. Notably, AutoAttack---a strong ensemble attack---yields the lowest robust accuracy at 94.8\%, confirming it as the most challenging perturbation method in our evaluation.

\subsection{Federated Learning}

Table~\ref{tab:federated} compares FedAvg~\cite{McMahan2017}, FedProx~\cite{Li2020FedProx}, and our FedGTD protocol across convergence speed, final accuracy, and Byzantine fault tolerance.

\begin{table}[!t]
\centering
\caption{Federated Learning Protocol Comparison ($K{=}10$ clients)}
\label{tab:federated}
\renewcommand{\arraystretch}{1.15}
\begin{tabular}{@{}lccc@{}}
\toprule
\textbf{Protocol} & \textbf{Rounds} & \textbf{Final Acc (\%)} & \textbf{Byz.\ Tol.\ (\%)} \\
\midrule
FedAvg~\cite{McMahan2017}   & 85  & 96.4 & 0 \\
FedProx~\cite{Li2020FedProx} & 72  & 97.1 & 10 \\
\textbf{FedGTD (Ours)}       & \textbf{58} & \textbf{97.9} & \textbf{30} \\
\bottomrule
\end{tabular}
\end{table}

FedGTD converges in 58 communication rounds---32\% fewer than FedAvg---while achieving 97.9\% accuracy. Its geometric trimmed-mean aggregation tolerates up to 30\% Byzantine clients, making it suitable for adversarial federated environments where compromised nodes may inject poisoned updates.

\subsection{Ablation Study}

To quantify each component's contribution, we conduct an ablation study by removing one module at a time from the full RobustIDPS pipeline.

\begin{table}[!t]
\centering
\caption{Ablation Study on CICIDS2017 (Accuracy / Macro-F1)}
\label{tab:ablation}
\renewcommand{\arraystretch}{1.15}
\begin{tabular}{@{}lcc@{}}
\toprule
\textbf{Configuration} & \textbf{Acc (\%)} & \textbf{F1 (\%)} \\
\midrule
Full RobustIDPS                    & \textbf{99.2} & \textbf{98.7} \\
\midrule
$-$ Adversarial Training           & 98.4 & 97.6 \\
$-$ Continual Learning (EWC)       & 97.8 & 96.9 \\
$-$ RL Response Agent              & 99.1 & 98.5 \\
$-$ Federated Aggregation          & 98.9 & 98.2 \\
$-$ Heterogeneous Ensemble (CNN only) & 98.0 & 97.1 \\
$-$ Attention Fusion               & 98.6 & 97.8 \\
\bottomrule
\end{tabular}
\end{table}

The heterogeneous ensemble and continual learning modules contribute the largest individual gains (1.2\% and 1.4\% accuracy, respectively). Removing adversarial training causes a 0.8\% accuracy drop under clean conditions but a far larger degradation under attack, as shown in Section~\ref{sec:results}D. The RL response agent has minimal effect on detection accuracy, as expected, since its role is downstream action selection rather than classification.

%% ============================================================
%% SECTION VII: DISCUSSION
%% ============================================================
\section{Discussion}
\label{sec:discussion}

\textbf{Limitations.} Although RobustIDPS demonstrates strong performance across multiple benchmarks, several limitations warrant discussion. First, adversarial training increases computational cost by approximately 2.8$\times$ relative to standard training, which may limit applicability in resource-constrained edge deployments. Second, continual learning with EWC introduces per-parameter Fisher information overhead that scales linearly with model size. Third, our evaluation relies on publicly available datasets that may not fully capture the diversity of production network traffic.

\textbf{Computational Cost.} End-to-end training of the full pipeline requires approximately 14 GPU-hours on four A100s. Inference latency averages 3.2\,ms per flow on a single GPU, meeting the sub-5\,ms requirement for inline deployment at 10\,Gbps link speeds. The federated protocol adds communication overhead proportional to the model size ($\sim$4.2\,MB per round for 10 clients).

\textbf{Deployment Considerations.} The system is currently deployed at \texttt{robustidps.ai} as a proof-of-concept platform. Production deployment would require integration with existing SIEM infrastructure, rate limiting of the RL response agent to prevent automated actions from cascading, and periodic retraining as the threat landscape evolves. The modular architecture facilitates selective deployment---organizations may adopt the detection ensemble without the federated component if data sovereignty constraints are satisfied locally.

\textbf{Ethical Considerations.} Automated intrusion response systems carry inherent risks of false positive-triggered disruption. Our CPO-based safety constraints mitigate but do not eliminate this risk. We recommend human-in-the-loop oversight for high-consequence response actions such as network isolation.

%% ============================================================
%% SECTION VIII: CONCLUSION
%% ============================================================
\section{Conclusion}
\label{sec:conclusion}

We presented RobustIDPS, a unified framework for network intrusion detection and prevention that integrates heterogeneous deep ensemble classification, EWC-based continual learning, constrained reinforcement learning for automated response, multi-attack adversarial training, and Byzantine-tolerant federated learning. Extensive experiments on five benchmark datasets demonstrate that RobustIDPS achieves state-of-the-art detection accuracy (99.2\% on CICIDS2017) while maintaining robust performance under six adversarial attack models, tolerating up to 30\% Byzantine clients in federated settings, and limiting catastrophic forgetting to under 1\% backward transfer across five sequential tasks.

Future work will explore transformer-based flow encoders, graph neural networks for lateral movement detection, and real-time adaptation mechanisms that combine continual and meta-learning for few-shot recognition of zero-day attacks.

%% ============================================================
%% BIBLIOGRAPHY
%% ============================================================
\balance

\begin{thebibliography}{99}

\bibitem{Kirkpatrick2017}
J.~Kirkpatrick \emph{et~al.}, ``Overcoming catastrophic forgetting in neural networks,'' \emph{Proc. Nat. Acad. Sci.}, vol.~114, no.~13, pp.~3521--3526, 2017.

\bibitem{Goodfellow2015}
I.~J.~Goodfellow, J.~Shlens, and C.~Szegedy, ``Explaining and harnessing adversarial examples,'' in \emph{Proc. ICLR}, 2015.

\bibitem{Madry2018}
A.~Madry, A.~Makelov, L.~Schmidt, D.~Tsipras, and A.~Vladu, ``Towards deep learning models resistant to adversarial attacks,'' in \emph{Proc. ICLR}, 2018.

\bibitem{Carlini2017}
N.~Carlini and D.~Wagner, ``Towards evaluating the robustness of neural networks,'' in \emph{Proc. IEEE Symp. Security and Privacy}, pp.~39--57, 2017.

\bibitem{McMahan2017}
B.~McMahan, E.~Moore, D.~Ramage, S.~Hampson, and B.~A.~y~Arcas, ``Communication-efficient learning of deep networks from decentralized data,'' in \emph{Proc. AISTATS}, pp.~1273--1282, 2017.

\bibitem{Schulman2015}
J.~Schulman, S.~Levine, P.~Abbeel, M.~Jordan, and P.~Moritz, ``Trust region policy optimization,'' in \emph{Proc. ICML}, pp.~1889--1897, 2015.

\bibitem{Achiam2017}
J.~Achiam, D.~Held, A.~Tamar, and P.~Abbeel, ``Constrained policy optimization,'' in \emph{Proc. ICML}, pp.~22--31, 2017.

\bibitem{Sharafaldin2018}
I.~Sharafaldin, A.~H.~Lashkari, and A.~A.~Ghorbani, ``Toward generating a new intrusion detection dataset and intrusion traffic characterization,'' in \emph{Proc. ICISSP}, pp.~108--116, 2018.

\bibitem{Moustafa2015}
N.~Moustafa and J.~Slay, ``UNSW-NB15: A comprehensive data set for network intrusion detection systems,'' in \emph{Proc. MilCIS}, pp.~1--6, 2015.

\bibitem{Tavallaee2009}
M.~Tavallaee, E.~Bagheri, W.~Lu, and A.~A.~Ghorbani, ``A detailed analysis of the KDD CUP 99 data set,'' in \emph{Proc. IEEE CISDA}, pp.~1--6, 2009.

\bibitem{Neto2023}
E.~C.~P.~Neto \emph{et~al.}, ``CIC-IoT-Dataset-2023: A dataset for IoT device activity detection and classification,'' \emph{Internet of Things}, vol.~24, 2023.

\bibitem{CSEIDS2018}
Communications Security Establishment (CSE) and Canadian Institute for Cybersecurity (CIC), ``CSE-CIC-IDS2018,'' 2018. [Online]. Available: \url{https://www.unb.ca/cic/datasets/ids-2018.html}

\bibitem{Schulman2017}
J.~Schulman, F.~Wolski, P.~Dhariwal, A.~Radford, and O.~Klimov, ``Proximal policy optimization algorithms,'' \emph{arXiv preprint arXiv:1707.06347}, 2017.

\bibitem{Kingma2015}
D.~P.~Kingma and J.~Ba, ``Adam: A method for stochastic optimization,'' in \emph{Proc. ICLR}, 2015.

\bibitem{Chen2016}
T.~Chen and C.~Guestrin, ``XGBoost: A scalable tree boosting system,'' in \emph{Proc. ACM KDD}, pp.~785--794, 2016.

\bibitem{Lopez-Paz2017}
D.~Lopez-Paz and M.~Ranzato, ``Gradient episodic memory for continual learning,'' in \emph{Proc. NeurIPS}, pp.~6467--6476, 2017.

\bibitem{Li2020FedProx}
T.~Li, A.~K.~Sahu, M.~Zaheer, M.~Sanjabi, A.~Talwalkar, and V.~Smith, ``Federated optimization in heterogeneous networks,'' in \emph{Proc. MLSys}, 2020.

\bibitem{Moosavi2016}
S.-M.~Moosavi-Dezfooli, A.~Fawzi, and P.~Frossard, ``DeepFool: A simple and accurate tool to fool deep neural networks,'' in \emph{Proc. IEEE CVPR}, pp.~2574--2582, 2016.

\bibitem{Croce2020}
F.~Croce and M.~Hein, ``Reliable evaluation of adversarial robustness with an ensemble of attacks,'' in \emph{Proc. ICML}, pp.~2206--2216, 2020.

\bibitem{Lin2022}
Z.~Lin, Y.~Shi, and Z.~Xue, ``IDSGAN: Generative adversarial networks for attack generation against intrusion detection,'' \emph{Pacific-Asia Conf. Knowl. Discovery Data Mining}, pp.~79--91, 2022.

\bibitem{Zenke2017}
F.~Zenke, B.~Poole, and S.~Ganguli, ``Continual learning through synaptic intelligence,'' in \emph{Proc. ICML}, pp.~3987--3995, 2017.

\bibitem{Shin2017}
H.~Shin, J.~K.~Lee, J.~Kim, and J.~Kim, ``Continual learning with deep generative replay,'' in \emph{Proc. NeurIPS}, pp.~2990--2999, 2017.

\bibitem{Yin2018}
D.~Yin, Y.~Chen, R.~Kannan, and P.~Bartlett, ``Byzantine-resilient distributed learning: Towards optimal statistical rates,'' in \emph{Proc. ICML}, pp.~5650--5659, 2018.

\bibitem{Blanchard2017}
P.~Blanchard, E.~M.~El~Mhamdi, R.~Guerraoui, and J.~Stainer, ``Machine learning with adversaries: Byzantine tolerant gradient descent,'' in \emph{Proc. NeurIPS}, pp.~119--129, 2017.

\bibitem{Mirsky2018}
Y.~Mirsky, T.~Doitshman, Y.~Elovici, and A.~Shabtai, ``Kitsune: An ensemble of autoencoders for online network intrusion detection,'' in \emph{Proc. NDSS}, 2018.

\bibitem{Vinayakumar2019}
R.~Vinayakumar, M.~Alazab, K.~P.~Soman, P.~Poornachandran, A.~Al-Nemrat, and S.~Venkatraman, ``Deep learning approach for intelligent intrusion detection system,'' \emph{IEEE Access}, vol.~7, pp.~41525--41550, 2019.

\bibitem{Ferrag2020}
M.~A.~Ferrag, L.~Maglaras, S.~Moschoyiannis, and H.~Janicke, ``Deep learning for cyber security intrusion detection: Approaches, datasets, and comparative study,'' \emph{J. Inf. Security Appl.}, vol.~50, 2020.

\bibitem{Pujari2022}
M.~Pujari, Y.~Pacheco, B.~Satam, and S.~Hariri, ``A comparative study on deep learning approaches for intrusion detection in IoT networks,'' \emph{IEEE Internet Things J.}, vol.~9, no.~12, pp.~9849--9860, 2022.

\bibitem{Karimipour2019}
H.~Karimipour, A.~Dehghantanha, R.~M.~Parizi, K.-K.~R.~Choo, and H.~Leung, ``A deep and scalable unsupervised machine learning system for cyber-attack detection in large-scale smart grids,'' \emph{IEEE Access}, vol.~7, pp.~80778--80788, 2019.

\bibitem{Liu2022}
H.~Liu, B.~Lang, M.~Liu, and H.~Yan, ``CNN and RNN based payload classification methods for attack detection,'' \emph{Knowl.-Based Syst.}, vol.~163, pp.~332--341, 2022.

\bibitem{Aburomman2017}
A.~A.~Aburomman and M.~B.~I.~Reaz, ``A novel SVM-kNN-PSO ensemble method for intrusion detection system,'' \emph{Appl. Soft Comput.}, vol.~38, pp.~360--372, 2017.

\bibitem{Khraisat2019}
A.~Khraisat, I.~Gondal, P.~Vamplew, and J.~Kamruzzaman, ``Survey of intrusion detection systems: Techniques, datasets and challenges,'' \emph{Cybersecurity}, vol.~2, no.~1, pp.~1--22, 2019.

\bibitem{Xu2018}
W.~Xu, D.~Evans, and Y.~Qi, ``Feature squeezing: Detecting adversarial examples in deep neural networks,'' in \emph{Proc. NDSS}, 2018.

\bibitem{Wang2020}
S.~Wang, T.~Tuor, T.~Salonidis, K.~K.~Leung, C.~Makaya, T.~He, and K.~Chan, ``Adaptive federated learning in resource constrained edge computing systems,'' \emph{IEEE J. Sel. Areas Commun.}, vol.~37, no.~6, pp.~1205--1221, 2020.

\bibitem{Yang2019}
Q.~Yang, Y.~Liu, T.~Chen, and Y.~Tong, ``Federated machine learning: Concept and applications,'' \emph{ACM Trans. Intell. Syst. Technol.}, vol.~10, no.~2, pp.~1--19, 2019.

\bibitem{Tesauro2007}
G.~Tesauro, ``Reinforcement learning in autonomic computing: A manifesto and case studies,'' \emph{IEEE Internet Comput.}, vol.~11, no.~1, pp.~22--30, 2007.

\end{thebibliography}

\end{document}
